\documentclass[a4paper,12pt]{article}
\usepackage[utf8x]{inputenc}
\usepackage[T1]{fontenc}
\usepackage[english]{babel}
\usepackage{lmodern}
\usepackage[pdftex,
            pdfauthor={Alexis Marchand},
            pdftitle={Mathematicians, machines, and the Earth},
            colorlinks=true]{hyperref}
\usepackage{filecontents}

\title{Mathematicians, machines, and the Earth}
\author{Alexis Marchand}

\begin{filecontents}{refs.bib}
	@misc{chu,
		author={Chu, Tasmin},
		title={The AI dissenter viewpoint},
		year={2026},
		url={https://proofsandprompts.com/2026/08/09/the-ai-dissenter-viewpoint/},
	}
	
	@misc{weinreich,
		author={Weinreich, Max},
		title={The crisis of AI-generated mathematics},
		year={2026},
		url={https://arxiv.org/abs/2608.02859},
	}
	
	
	@misc{unu,
		author={M. Aczel and S. Chamanara and M. Matin and A. Farsi and T. Marwala and K. Madani},
		title={Environmental Cost of AI's Energy Use: Carbon, Water and Land Footprints},
		year={2026},
		note={UNU-INWEH Report, United Nations University Institute for Water, Environment and Health (UNU-INWEH), Richmond Hill, Ontario},
		url={https://doi.org/10.53328/INR26RMA002}
	 }
	
	@article{ding-etal,
		title = {Tracking the carbon footprint of global generative artificial intelligence},
		journal = {The Innovation},
		volume = {6},
		number = {5},
		pages = {100866},
		year = {2025},
		issn = {2666-6758},
		author = {Zhaohao Ding and Jianxiao Wang and Yiyang Song and Xiaokang Zheng and Guannan He and Xiupeng Chen and Tiance Zhang and Wei-Jen Lee and Jie Song},
		url = {https://doi.org/10.1016/j.xinn.2025.100866}
	}

	@inbook{lambert-luccioni,
		author = {Lambert, Katherine and Luccioni, Sasha},
		title = {From Cradle to Cloud: A Life Cycle Review of AI's Environmental Footprint},
		year = {2026},
		isbn = {9798400725968},
		publisher = {Association for Computing Machinery},
		address = {New York, NY, USA},
		booktitle = {Proceedings of the 2026 ACM Conference on Fairness, Accountability, and Transparency},
		pages = {518–540},
		numpages = {23},
		url = {https://doi.org/10.1145/3805689.3812401},
	}
\end{filecontents}

\usepackage[backend=bibtex,style=nature,sorting=nyt]{biblatex}
\addbibresource{refs.bib}

\begin{document}
	\maketitle
	
	In the summer of 2026, mathematicians are watching in shock as, conjecture after conjecture, large language models\footnote{I intentionally avoid the term `artificial intelligence' as I believe it is helping the marketing strategy of the companies developing those tools, playing on a belief that their products are more than merely machines.} are proving their ability to write research papers solving important mathematical problems.
	In the same summer, Europeans are watching in shock as heatwaves, droughts, and wildfires are following each other with a yet unseen intensity in Western Europe --- and similarly extreme climatic events have occured around the globe in the past few years.
	
	Two crises are unfolding at the same time, a \emph{climate crisis} and a \emph{large language model crisis}.
	I believe that those two crises are profoundly intertwined; they feed each other and have the same deep roots, and they should be thought together.
	Those roots are a Promethean belief in the might of human technology and a thirst for everlasting growth.
	These led to the Industrial Revolution in the nineteenth century, with increases in productivity that improved our life conditions, but are now confronted with the resource limitations of our planet, and have modified our climate in a destructive manner.
	Today, the development of large language models is coming with similar promises of infinite growth and unlimited power, while relying on resource-intensive data centres\footnote{As an example, according to a UNU report \cite{unu}, LLM workloads were responsible for about 20\% of the global electricity consumption of data centres in 2025, or about 90 TWh, and by 2030, the annual electricity consumption of large language models could reach about 380 TWh, which is comparable with the total electricity consumption of France (470 TWh) or Saudi Arabia (420 TWh) in 2025. See also another study by Ding et al. \cite{ding-etal} and a literature review by Lambert and Luccioni \cite{lambert-luccioni}.}, as well as profit-driven, wealth-accumulating leadership.
	This is aggravating the current climate crisis and is relying on the same way of thinking that has shown its limitations.
	It is also worth noting that the Industrial Revolution facilitated the development of the mass destruction weapons that were the instruments of two traumatic world wars; in a time of geopolitical tensions --- that are not unrelated to the resource limitations of the Earth --- we should carefully think about what the new large language model revolution might lead to.
	
	\paragraph{The role of mathematicians.}
	Where those two crises differ is in the involvement of mathematicians.
	With respect to the climate crisis, I believe that mathematicians are not different from other human beings: we need to act, but just in the same way that all of society needs to act.
	Before going any further, let me use one more sentence to emphasize: \emph{we do need to solve the climate crisis}; we need to rethink human activity in order to reduce our carbon emissions; as mathematicians, we need in particular to question the very carbon-intensive manner in which we tend to travel, as well as a certain pressure for productivity which pushes us to write more and more papers.
	
	But with respect to the large language model crisis which is promising to worsen the already existing climate crisis, I believe that mathematicians do play a central role, in three ways.
	\begin{enumerate}
		\item Large language models are developed using \emph{mathematics}.
		\item The creators of large language models are mostly \emph{mathematicians} or at least have a mathematical training, both within and outside academia.
		\item Large language model companies are using \emph{mathematical research} as the spearhead of their marketing strategy, hoping to convince the public that a machine that can solve conjectures about which mathematicians have thought for decades must be intelligent.
	\end{enumerate}
	The involvement of mathematics at the core of large language models gives us, as mathematicians, a societal responsibility that is probably unprecedented in the history of our field; we have a duty to lead the way to the human response to the looming crisis.
	
	\paragraph{Acting as a community of mathematicians.}
	In response to the involvement of mathematics in large language models, we do need to act.
	\begin{enumerate}
		\item Since large language models rely on mathematics, we have a duty to explain their inner working to the public.
		We must demystify them and insist that they are not mysterious intelligent beings, but tools whose mechanism can be understood and controlled.
		\item The present architects of large language models are our current or former colleagues, while the future ones are our students; we need to maintain a dialogue with them, discussing the dangers of those tools, the philosophy that is driving their development, and pushing for the technology to move towards virtuous applications.
		\item We must question and discuss the reasons why we are doing mathematics.
		This is already happening and a lot of beautiful ideas are emerging from the ongoing debates.
		We urgently need to communicate those reasons to the public, convey the beauty of mathematics, and explain why we cannot be replaced by paper-writing machines.
		We also need to stop participating in the buzz around mathematics created by large language models, and stop giving benchmark problems to the companies behind those models.
		I think that, until now, we have not been very good at explaining the motivation for mathematics: who among us has never been asked: ``If there is no application, then what is the point?''
		We must convince society that it is worthwile for human beings to pursue mathematics; our jobs depend on it.
	\end{enumerate}
	All of the above points push us to spend more energy \emph{communicating} mathematics outside the world of mathematicians.
	We must get out of our ivory tower and exchange with the world; such communication activities should become a full part of our job, rather than a side hobby as they are sometimes considered today.
	
	\paragraph{Acting as individual mathematicians.}
	In view of the contribution of large language models to the climate crisis, we need to question our use of them.
	We need to do so before our lifestyles are altered in a way that might make it difficult to go back, similar to what has happened with the Industrial Revolution and the technologies that came out of it.
	\begin{itemize}
		\item The use of large language models to create and correct language is becoming increasingly widespread.
		While this is not specific to mathematics, communicating ideas is an essential part of our work, so I think it is worth taking a moment to think about this.
		Human beings do not have a perfect mastery of language: they make mistakes, and they find the process of writing slow and painful.
		As for the first point, I would like to ask: is it really worth consuming vast amounts of resources to remove a few typos from a paper or email?
		I personally think that I can live with the certainty that my writing will keep containing mistakes as it always has.
		As for the second point: yes, producing language is difficult, but I believe that it is at the core of human thoughts and society; if we stop practising language creation, then we might become unable to develop and communicate complex ideas.
		Do we really want to \emph{delegate language to machines}?
		This is also true in mathematics, and the communication of ideas through language is inseparable from the creative process in mathematics.
		\item The second, more and more prevalent manner, in which large language models are being used in mathematics is for the discovery and creation of mathematical ideas.
		The deleterious effects of such practices are being discussed at length elsewhere, so let me be brief and point to the excellent essays of Tasmin Chu \cite{chu} and Max Weinreich \cite{weinreich}.
		Large language models risk leading to a regime where mathematics is produced extremely fast, with the community unable to keep track, and no expert able to understand and explain new ideas.
		Mathematics will lose its meaning if it does not remain a human activity.
	\end{itemize}
	
	Considering the role of large language models in aggravating the climate crisis, considering the harm they can produce in mathematics and human thought, and considering our responsibility as mathematicians, my personal choice is the following.
	
	\medskip
	
	\emph{I will not (knowingly) use any large language model in my work and private life.}
	
	\medskip
	
	I will refuse to review papers making (significant) use of large language models, I will refuse to collaborate on projects using large language models, and I will not consent to any of my writing being fed to large language models.
	I will keep to this stance as long as large language models are based on the Earth- and human-damaging philosophy which is currently at their core.
	
	I understand that this position might sound radical, and I do not expect all mathematicians to follow it, but my hope in writing this text is to encourage everyone to think about the impact of our choices, not just on mathematics, but on the Earth and its inhabitants.
	
	\printbibliography

\end{document}